% The HTML version of the manual is produced by lwarp; see the Makefile.
% Macro \CFmultipage, if defined, selects multi-page HTML output (one page per
% chapter).  Otherwise, the HTML output is a single page, manual.html.

% Text fonts and encodings must be loaded before lwarp.
\usepackage[T1]{fontenc}
\usepackage[utf8]{inputenc}
\usepackage{pslatex}

\ifdefined\CFmultipage
\usepackage[latexmk,HomeHTMLFilename=manual-multipage,HTMLFilename={manual-}]{lwarp}
\setcounter{FileDepth}{0}
\else
\usepackage[latexmk,HomeHTMLFilename=manual,HTMLFilename={}]{lwarp}
\setcounter{FileDepth}{-5}
\fi
\CSSFilename{manual.css}
% Workarounds for lwarp problems, used by the Checker Framework manual and by
% the Annotation File Utilities manual.  Input this file after loading lwarp.
% It affects only the HTML version.

\begin{warpHTML}

% lwarp escapes < within Verbatim, but not in ordinary text, \texttt, or
% alltt.  A browser would treat text such as "List<Object>" as an HTML tag and
% not display it.  So, within the document, < is an active character that
% produces an HTML entity.  (An unescaped > is legal in HTML text.  Leaving >
% alone lets it continue to delimit arguments, as in the manual's \<...>.)
\begingroup
\catcode`\<=\active
\gdef<{\ifmmode\string<\else\textless\fi}
\endgroup
\AfterEndPreamble{\catcode`\<=\active}
% \@noligs, which alltt calls, redefines < to print itself.
% Ligatures are irrelevant in HTML.
\makeatletter
\def\verbatim@nolig@list{\do\`\do\>\do\,\do\'\do\-}
\makeatother

% BibTeX's alpha style defines \etalchar using math mode, which lwarp would
% convert to an image, and which fails within a citation.
\newcommand{\etalchar}[1]{\textsuperscript{#1}}
% The .bbl file uses \newcommand to define \etalchar.
\pretocmd{\bibliography}{\let\etalchar\relax}{}{}
\AtBeginEnvironment{thebibliography}{\renewcommand{\etalchar}[1]{\textsuperscript{#1}}}

\end{warpHTML}

\HTMLTitle{The Checker Framework Manual}
\HTMLDescription{The Checker Framework Manual: Custom pluggable types for Java}

\usepackage{microtype}
\usepackage{fullpage}
\usepackage{graphicx}
\usepackage{url}
\usepackage{color}
% Not supported by lwarp: \usepackage{minitoc}
\usepackage{proof}

\usepackage{relsize}
% \def\codesize{\smaller}
\def\codesize{\relax}           % for "pslatex"

\begin{warpprint}
\newcommand{\code}[1]{\ifmmode{\mbox{\codesize\ttfamily{#1}}}\else{\codesize\ttfamily #1}\fi}
%\DeclareRobustCommand{\code}[1]{\codesize{\path{#1}}}
\end{warpprint}
\begin{warpHTML}
\newcommand{\code}[1]{\texttt{#1}}
\end{warpHTML}
\def\<#1>{\code{#1}}

% For an internal link to a LaTeX label, use \hyperref[label]{anchor text}.

% This can't handle a URL with an embedded "#" -- at least at UW CSE
\newcommand{\myurl}[1]{{\codesize\url{#1}}}
% A command that shows the URL in the printed manual.
% Make a URL visible in the PDF but just be attached to anchor text in HTML:
\newcommand{\ahreforurl}[2]{%
\ifbool{warpingHTML}{\href{#1}{#2}}{#2 (\url{#1})}}

\usepackage{listings}
\usepackage{alltt}
\usepackage{fancyvrb}
\begin{warpprint}
\RecustomVerbatimEnvironment{Verbatim}{Verbatim}{fontsize=\codesize}
\end{warpprint}

% In HTML, \smaller has no effect, so the BlockClass sets the font size.
\newenvironment{mysmall}
  {\begin{BlockClass}[font-size:small]{mysmall}\smaller}
  {\end{BlockClass}}
\newenvironment{myxsmall}
  {\begin{BlockClass}[font-size:x-small]{myxsmall}\smaller\smaller}
  {\end{BlockClass}}
\newenvironment{myxxsmall}
  {\begin{BlockClass}[font-size:xx-small]{myxxsmall}\smaller\smaller\smaller}
  {\end{BlockClass}}

% The argument is the HEIGHT, not the width.
\newcommand{\includeimage}[2]{%
\begin{center}
\includeimagenocentering{#1}{#2}%
\warpprintonly{\vspace{-1.5\baselineskip}}%
\end{center}}

% The argument is the HEIGHT, not the width.  It is used only in print.
% In HTML, the image appears at its natural size (the .svg file is used).
\newcommand{\includeimagenocentering}[2]{%
\ifbool{warpingHTML}%
{\includegraphics{figures/#1}}%
{\resizebox{!}{#2}{\includegraphics{figures/#1}}}}


% At least 80% of every float page must be taken up by
% floats; there will be no page with more than 20% white space.
\def\topfraction{.8}
\def\dbltopfraction{\topfraction}
\def\floatpagefraction{\topfraction}     % default .5
\def\dblfloatpagefraction{\topfraction}  % default .5
\def\textfraction{.2}


% A symbol that is typeset as math in print, and as a Unicode character in
% HTML, where lwarp would otherwise render the math as an image.
% Arg 1: hexadecimal Unicode code point.  Arg 2: math-mode command.
\newcommand{\textsym}[2]{\ifbool{warpingHTML}{\HTMLunicode{#1}}{\ensuremath{#2}}}
\newcommand{\Rightarrowsym}{\textsym{21D2}{\Rightarrow}}
\newcommand{\rightarrowsym}{\textsym{2192}{\rightarrow}}
\newcommand{\leftrightarrowsym}{\textsym{2194}{\leftrightarrow}}
\newcommand{\precsym}{\textsym{227A}{\prec}}
\newcommand{\capsym}{\textsym{2229}{\cap}}
\newcommand{\neqsym}{\textsym{2260}{\neq}}
\newcommand{\lesym}{\textsym{2264}{\le}}
\newcommand{\timessym}{\textsym{00D7}{\times}}
\newcommand{\approxsym}{\textsym{2248}{\approx}}

% Left and right curly braces and backslash, in tt font
\newcommand{\ttlcb}{\texttt{\char "7B}}
\newcommand{\ttrcb}{\texttt{\char "7D}}
\newcommand{\ttbs}{\texttt{\char "5C}}


\begin{warpprint}
  %% Bring items closer together in list environments
  % Prevent infinite loops
  \let\Itemize =\itemize
  \let\Enumerate =\enumerate
  \let\Description =\description
  % Zero the vertical spacing parameters
  \def\Nospacing{\itemsep=0pt\topsep=0pt\partopsep=0pt\parskip=0pt\parsep=0pt}
  % Redefine the environments in terms of the original values
  \renewenvironment{itemize}{\Itemize\Nospacing}{\endlist}
  \renewenvironment{enumerate}{\Enumerate\Nospacing}{\endlist}
  \renewenvironment{description}{\Description\Nospacing}{\endlist}

  % Add line between figure and text
  \makeatletter
  \def\topfigrule{\kern3\p@ \hrule \kern -3.4\p@} % the \hrule is .4pt high
  \def\botfigrule{\kern-3\p@ \hrule \kern 2.6\p@} % the \hrule is .4pt high
  \def\dblfigrule{\kern3\p@ \hrule \kern -3.4\p@} % the \hrule is .4pt high
  \makeatother
\end{warpprint}


% Reference to Checker Framework Javadoc for a class (not a method, etc.).
% Arg 1: directory under org/checkerframework/, including internal "/", but
% no leading or trailing "/".
% Arg 2: class name.
% In the printed version, only the base class name appears.
% In the HTML version, it's a link to the Javadoc.
\newcommand{\refclass}[2]{\href{../api/org/checkerframework/#1/#2.html}{\<#2>}}

% Reference to Checker Framework Javadoc for a type qualifier/annotation class
% (not a method, etc.).  Like \refclass, but prepends an @ to the qualifier name.
\newcommand{\refqualclass}[2]{\href{../api/org/checkerframework/#1/#2.html}{\<@#2>}}

% Reference to Checker Framework Javadoc for a type qualifier/annotation class
% (not a method, etc.) that takes parameters.  This prepends an @ to the
% qualifier name and includes parentheses around the parameters.
\newcommand{\refqualclasswithparams}[3]{\href{../api/org/checkerframework/#1/#2.html}{\<@#2>}\<(\allowbreak #3)>}

% Reference to Checker Framework Javadoc for a method or field.
% Arg 1: package name under org/checkerframework, using "/" as a separator,
% with no leading or trailing "/". example: "checker/nullness/qual"
% Arg 2: class name.
% Arg 3: method name.
% Arg 4: fully-qualified arguments.  Example: "(T)".  For Java 8, may need
% to use "-T-" with dashes instead of parentheses.
% In the printed version, only "class.method" appears.
% In the HTML version, it's a link to the Javadoc.
\newcommand{\refmethod}[4]{\href{../api/org/checkerframework/#1/#2.html\##3#4}{\<#2.#3>}}
% Omits the class name in the visible text; is terser.  Good for mentioning
% a method that needs to be overridden.
\newcommand{\refmethodterse}[4]{\href{../api/org/checkerframework/#1/#2.html\##3#4}{\<#3>}}
% Permits specification of the anchor text.  Good for mentioning the
% arguments of an overloaded method that needs to be overridden.
\newcommand{\refmethodanchortext}[5]{\href{../api/org/checkerframework/#1/#2.html\##3#4}{\<#5>}}
% Like refmethod, but formatted more concisely as ``field'' instead of
% ``Class.field''.  The last argument is usually empty {}.
\newcommand{\refenum}[4]{\href{../api/org/checkerframework/#1/#2.html\##3#4}{\<#3>}}

% Reference to Oracle (previously Sun) Javadoc.
% Arg 1: .html reference, without the .../api/ prefix.  Quote "#" as "\#" if it appears.
% Arg 2: What will appear in the formatted manual.
% Do not use \url in the body of any macro.
\newcommand{\sunjavadoc}[2]{\href{https://docs.oracle.com/en/java/javase/25/docs/api/#1}{\<#2>}}
% Like \sunjavadoc, but for annotations and prepends "@".
\newcommand{\sunjavadocanno}[2]{\href{https://docs.oracle.com/en/java/javase/25/docs/api/#1}{\<@#2>}}
% Now, for Java EE (version 7 is the last version):
\newcommand{\javaeejavadoc}[2]{\href{https://docs.oracle.com/javaee/7/api/#1}{\<#2>}}
\newcommand{\javaeejavadocanno}[2]{\href{https://docs.oracle.com/javaee/7/api/#1}{\<@#2>}}


% These commands are for long-range references:  those not in the same chapter.
% In HTML, the page number is omitted.
\newcommand{\refwithpage}[1]{\ref{#1}\warpprintonly{, page~\pageref{#1}}}
\newcommand{\refwithpageparen}[1]{\ref{#1}\warpprintonly{ (page~\pageref{#1})}}
% Use \chapterpageref to reference only chapters, not sections.
\newcommand{\chapterpageref}[1]{Chapter~\refwithpage{#1}}
\newcommand{\sectionpageref}[1]{Section~\refwithpage{#1}}


% In HTML, each sectioning command is followed by a link to itself, so that
% users can easily obtain a URL for the section.  The optional argument keeps
% the link icon out of the table of contents and out of HTML file names.
\ifbool{warpingHTML}{%
% The link icon is an image, figures/chainlink.svg, specified in manual.css.
% (lwarp drops \includegraphics within a sectioning command.)
\newcommand{\selflink}[1]{\ \hyperref[#1]{\InlineClass{selflink}{}}}
}{%
\newcommand{\selflink}[1]{}
}
\newcommand{\chapterAndLabel}[2]{\chapter[{#1}]{#1\selflink{#2}}\label{#2}}
\newcommand{\sectionAndLabel}[2]{\section[{#1}]{#1\selflink{#2}}\label{#2}}
\newcommand{\subsectionAndLabel}[2]{\subsection[{#1}]{#1\selflink{#2}}\label{#2}}
\newcommand{\subsubsectionAndLabel}[2]{\subsubsection[{#1}]{#1\selflink{#2}}\label{#2}}
\newcommand{\paragraphAndLabel}[2]{\paragraph[{#1}]{#1\selflink{#2}}\label{#2}}

% hyperref should be loaded after other packages.
\usepackage{hyperref}

%  LocalWords:  fontsize mysmall myxsmall myxxsmall subsubsection
